\chapter{Polytrauma Imaging}

\label{ch:polytrauma}

Polytrauma refers to multiple, severe, life-threatening injuries affecting more than one anatomical region. In the initial "golden hour" of trauma resuscitation, rapid and systematic imaging is crucial to identify hidden internal hemorrhages, organ lacerations, and skeletal instability. This chapter details the polytrauma imaging algorithm, the role of whole-body CT (WBCT), and the classic radiological signs of life-threatening trauma.

---

\section*{1. The Trauma Imaging Algorithm}
Imaging in polytrauma follows a strict algorithm based on the patient's hemodynamic stability:

\subsection*{Bedside Evaluation (Hemodynamically Unstable Patients)}
For patients who remain hemodynamically unstable, bedside assessment and immediate hemorrhage control take priority; transfer to CT is appropriate only when the treating trauma team judges it safe and useful. Bedside imaging may include:
\textbullet \textbf{eFAST (Extended Focused Assessment with Sonography for Trauma)}: Done immediately at the bedside to check for free fluid/blood (Morison's pouch, spleno-renal, pelvic pouch, and pericardium) and thoracic pneumothorax/hemothorax.
\textbullet \textbf{Portable Chest X-Ray}: Evaluates endotracheal tube placement, checks for tension pneumothorax, massive hemothorax, or wide mediastinum (suggestive of aortic injury).
\textbullet \textbf{Portable Pelvic X-Ray}: Identifies unstable pelvic ring fractures (e.g., "open-book" fractures) which can cause massive, fatal retroperitoneal bleeding.
\textbullet *Action*: In an unstable patient, a positive eFAST can support urgent hemorrhage control, but the destination and procedure depend on the suspected injury, physiology, and trauma-team assessment.

\subsection*{Whole-Body CT / "Pan-Scan" (Hemodynamically Stable Patients)}
For selected stable or sufficiently stabilized patients with suspected major multisystem injury, whole-body CT can rapidly define injury extent. It is not required for every trauma patient; use local trauma criteria and tailor coverage and phases to the suspected injuries. A typical examination may include:
1. \textbf{NCCT Head \& Cervical Spine}: Non-contrast scan to check for intracranial bleeding and cervical spinal fractures.
2. \textbf{CECT Chest (trauma protocol)}: Evaluates thoracic structures, including the aorta and active bleeding; a dedicated pulmonary angiographic phase is used only when indicated.
3. \textbf{CECT Abdomen \& Pelvis (Portal Venous Phase)}: Evaluates solid organs (liver, spleen, kidneys), bowel wall integrity, and mesenteric vascular injury.

---

\section*{2. Life-Threatening Traumatic Findings}

\subsection*{Intracranial Injuries}
Look for acute epidural or subdural hematomas, traumatic subarachnoid hemorrhage, mass effect, or midline shift. Identify signs of herniation (uncal herniation or tonsillar herniation).

\subsection*{Cervical Spine Fractures}
Trace the three spinal lines (anterior spinal, posterior spinal, and spinolaminar lines) on sagittal views. Any disruption represents instability. Unstable fractures include:
\textbullet \textbf{Jefferson Fracture}: Burst fracture of C1.
\textbullet \textbf{Hangman's Fracture}: Bilateral pedicle fractures of C2.

\subsection*{Thoracic Injuries}
\textbullet \textbf{Tension Pneumothorax}: Air in the pleural cavity compressing the lung and pushing the mediastinum and trachea toward the opposite side.
\textbullet \textbf{Aortic Injury / Transection}: Usually occurs at the \textbf{aortic isthmus} (just distal to the origin of the left subclavian artery). Look for a mediastinal hematoma and a focal aortic flap or pseudoaneurysm.
\textbullet \textbf{Pulmonary Contusion}: Ground-glass opacities or consolidation on CT that do not respect lobar borders, representing alveolar hemorrhage.

\subsection*{Abdominal Solid Organ Lacerations}
\textbullet \textbf{Spleen}: The most commonly injured solid organ in blunt abdominal trauma. Lacerations appear as linear hypodense areas within the enhancing spleen. Look for a subcapsular hematoma or active contrast extravasation ("blush"), indicating active arterial bleeding that requires embolization.
\textbullet \textbf{Liver}: Rated using the AAST (American Association for the Surgery of Trauma) grading scale. Look for periportal tracking of fluid and subcapsular hematomas.

\subsection*{Unstable Pelvic Fractures}
Disruption of the sacroiliac joints and pubic symphysis ("open-book" fracture) expands the pelvic volume, acting as a low-resistance space that can hold several liters of blood.

---

\section*{3. Structured Reporting Example: Positive WBCT in Polytrauma}

\begin{testbox}[Indications:]
34-year-old male involved in a high-speed motor vehicle collision. Intubated in the field. Hemodynamically stable but presenting with abdominal distension and pelvic instability.
\end{testbox}
\begin{refbox}[Findings:]
\textbf{Head \& Neck}: No intracranial hemorrhage or midline shift. Cervical spine alignment is normal. No fractures detected.\\
\textbf{Chest}: Endotracheal tube tip is positioned 4 cm above the carina. Bilateral pulmonary contusions are noted in the lower lobes. No pneumothorax or pleural effusion. The thoracic aorta is intact; no mediastinal hematoma.\\
\textbf{Abdomen \& Pelvis}:
\textbullet \textbf{Spleen}: There is a deep, irregular hypodensity traversing the splenic parenchyma, measuring 4.5 cm in length, consistent with a Grade III splenic laceration. No active contrast extravasation ("blush") is seen.
\textbullet \textbf{Liver, Pancreas, \& Kidneys}: Normal in appearance.
\textbullet \textbf{Peritoneum}: Moderate free fluid (hemoperitoneum) is noted in Morison's pouch and the pelvis.
\textbullet \textbf{Bones}: There is diastasis of the pubic symphysis measuring 3.2 cm, with disruption of the anterior sacroiliac ligaments. An associated large retroperitoneal hematoma is visualized in the pelvis, measuring 8 x 6 cm.
\end{refbox}
\begin{pearlbox}[Impression:]
1. Grade III splenic laceration with moderate hemoperitoneum. No active arterial contrast extravasation.
2. Unstable "open-book" pelvic fracture (pubic symphysis diastasis) with an associated pelvic retroperitoneal hematoma.
3. Bilateral pulmonary contusions.
\end{pearlbox}

\clearpage
